
\documentclass[10pt,twocolumn]{article}

\usepackage[margin=0.72in]{geometry}
\usepackage[T1]{fontenc}
\usepackage[utf8]{inputenc}
\usepackage{lmodern}
\usepackage{microtype}
\usepackage{amsmath,amssymb,amsthm,mathtools}
\usepackage{booktabs,tabularx,array,multirow}
\usepackage{graphicx}
\usepackage{subcaption}
\usepackage{xcolor}
\usepackage{hyperref}
\usepackage[nameinlink,capitalise,noabbrev]{cleveref}
\usepackage{enumitem}
\usepackage{siunitx}
\usepackage{float}
\usepackage{titlesec}
\usepackage{fancyhdr}
\usepackage{url}
\usepackage{csquotes}

\hypersetup{
  colorlinks=true,
  linkcolor=blue!55!black,
  citecolor=blue!55!black,
  urlcolor=blue!55!black,
  pdftitle={Text Detox Steering in Pretrained Representation Spaces},
  pdfauthor={Ilyas Galiev}
}

\setlength{\columnsep}{0.24in}
\setlength{\parindent}{0.9em}
\setlength{\parskip}{0pt}
\titleformat{\section}{\large\bfseries}{\thesection.}{0.45em}{}
\titleformat{\subsection}{\normalsize\bfseries}{\thesubsection.}{0.4em}{}
\titleformat{\subsubsection}{\normalsize\itshape}{\thesubsubsection.}{0.4em}{}

\pagestyle{fancy}
\fancyhf{}
\fancyhead[L]{Text Detox Steering in Pretrained Representation Spaces}
\fancyhead[R]{Galiev}
\fancyfoot[C]{\thepage}

\title{\vspace{-1.2em}\bfseries Safe Steering or Near-Identity Rewriting?\\
\vspace{0.25em}\large A Case Study of Text Detox in Pretrained Representation Spaces}
\author{Ilyas Galiev\\
\small Research note prepared from the Week 12 text-detox steering experiment\\
\small\href{https://github.com/Ily17as/Controllable-Latent-Steering-in-VAEs}{github.com/Ily17as/Controllable-Latent-Steering-in-VAEs}}
\date{April 2026}

\begin{document}
\maketitle

\begin{abstract}
This paper studies whether a human-interpretable text attribute---toxicity---can be controlled by linear steering directions in pretrained representation spaces. The pipeline embeds text with a sentence encoder, learns global and cluster-local detox directions on Civil Comments, rewrites toxic inputs with a text-to-text model, and compares fixed-step steering with automatic per-sample step selection on Civil Comments and ParaDetox evaluation subsets. The empirical picture is mixed but informative. Fixed steering produces substantially larger toxicity reduction than automatic selection, especially on ParaDetox, but at a noticeable semantic cost. In contrast, automatic step selection is highly conservative: the selected mean step size is effectively zero for the global method and almost zero for the local one, while feasible acceptance rates remain low (roughly 1--4\%). Local directions sometimes improve detox relative to global ones, especially on ParaDetox, but the effect is not decisive. The strongest result is therefore negative-but-useful: in the current setup, the dominant bottleneck is not the absence of locality alone, but the combination of candidate generation, feasibility constraints, and a selection rule that often prefers near-identity outputs.
\end{abstract}

\section{Introduction}
Controllable editing methods are useful only if they change the intended attribute without destroying the underlying content. In images, this question is often phrased as a trade-off between target gain and identity preservation. In text, an analogous setting is \emph{detoxification}: can we reduce toxicity while keeping meaning and discourse intent intact?

This note treats toxicity as a human-interpretable semantic factor and asks whether it can be controlled through steering directions in a pretrained representation space. The experiment is intentionally simple. A sentence encoder produces text embeddings, global and local directions are learned from clean-versus-toxic examples, and a separate rewriter converts steered prompts into candidate detoxified sentences. This setup isolates a useful scientific question:
\begin{quote}
Can representation steering produce controllable text detox, and if not, where does the pipeline fail?
\end{quote}

The answer is nuanced. The experiment does find a real detox signal: fixed steering reduces toxicity substantially on both evaluation sets. However, the most striking effect is not a clean victory of local over global steering. Instead, the strongest finding is that automatic per-sample step selection almost collapses to $\alpha^\star \approx 0$, while fixed steering trades stronger detox for weaker semantic preservation. This shifts the interpretation of the study. The main bottleneck is not simply \enquote{we need more local directions}; it is the interaction between representation quality, candidate generation, and the current selection objective.

The paper makes three concrete contributions:
\begin{enumerate}[leftmargin=1.2em,itemsep=0.15em]
\item a compact benchmark comparing global versus cluster-local detox directions and fixed versus automatic step selection;
\item an empirical diagnosis showing that the current automatic policy is extremely conservative and often selects near-identity outputs;
\item a failure-oriented interpretation of where the present text steering pipeline breaks and what should be strengthened next.
\end{enumerate}

\section{Problem setting and method}
\subsection{Representation steering for text detox}
Let $E(x)\in\mathbb{R}^d$ denote the embedding of an input text $x$. A steered representation is
\begin{equation}
z' = E(x) + \alpha\, d(x),
\end{equation}
where $d(x)$ is either a global detox direction $d_{\mathrm{glob}}$ or a cluster-local direction $d_{\mathrm{loc}}(x)$, and $\alpha$ is a steering magnitude selected from a small grid. The final output text is not decoded directly from the embedding. Instead, the steered representation conditions a text-to-text rewriting stage, which produces a bank of candidate rewrites.

The global direction is learned in two ways from Civil Comments embeddings: a normalized clean-minus-toxic centroid difference and a logistic-probe direction. The main experiment uses the probe direction. If the probe has coefficient vector $w$, then the detox direction is the normalized negative coefficient,
\begin{equation}
d_{\mathrm{probe}} = -\, \frac{w}{\|w\|_2},
\end{equation}
so that moving along $d_{\mathrm{probe}}$ shifts representations toward the cleaner side of the classifier boundary.

For local steering, the training embeddings are partitioned by MiniBatch K-Means into $K=8$ clusters. A separate detox direction is fit inside each cluster whenever both clean and toxic classes are sufficiently represented. In the archived run, all eight clusters satisfy the minimum class-count requirement, so there is no fallback to the global direction.

\subsection{Candidate generation and selection}
For each toxic input, the system generates rewrite candidates for a grid
\[
\alpha\in\{0, 0.5, 1.0, 1.5, 2.0\}.
\]
The candidate bank includes the original text and multiple rewritten outputs. Each candidate is scored by toxicity, semantic similarity to the original, and a surface-form penalty. Let $\Delta_{\mathrm{tox}}$ denote toxicity reduction, $\mathrm{sim}$ semantic similarity, and $\mathrm{surf}$ the surface penalty. The archived notebook uses the objective
\begin{equation}
J = w_t \,\mathrm{target\_sim}
    + w_d \,\Delta_{\mathrm{tox}}
    + w_s \,\mathrm{sim}
    - w_p \,\mathrm{surf},
\end{equation}
with weights $(w_t,w_d,w_s,w_p)=(0.35,0.35,0.25,0.05)$ and no fluency term in the final run.

A candidate is declared feasible if
\begin{equation}
\mathrm{sim}\ge 0.72
\quad\text{and}\quad
\Delta_{\mathrm{tox}}\ge 0.12.
\end{equation}
Fixed steering always reports the row at $\alpha=1.0$. Automatic selection chooses the highest-objective feasible row if one exists, and otherwise the highest-objective row overall.

\subsection{Data and models}
The steering directions are learned from a filtered Civil Comments subset. After thresholding clean and toxic examples and applying text-length filters, the archived run uses $9{,}942$ training rows. Evaluation uses $300$ toxic Civil Comments examples and $150$ ParaDetox examples. The sentence encoder is \texttt{all-mpnet-base-v2}, toxicity is scored by \texttt{unitary/toxic-bert}, and the rewriter is \texttt{google/flan-t5-base}. The global probe reaches training AUC $0.8931$ on Civil Comments embeddings.

\section{Experimental protocol}
The experiment compares four configurations on each evaluation set:
\begin{itemize}[leftmargin=1.2em,itemsep=0.1em]
\item global direction + automatic $\alpha^\star$,
\item global direction + fixed $\alpha=1.0$,
\item local direction + automatic $\alpha^\star$,
\item local direction + fixed $\alpha=1.0$.
\end{itemize}

The main reported metrics are mean toxicity drop, mean semantic similarity, feasible acceptance rate, and mean selected $\alpha$. Toxicity drop is defined as the difference between toxicity before and after rewriting, so larger is better. Semantic similarity is computed in the encoder space relative to the input. Acceptance rate is the fraction of selected outputs that satisfy the feasibility constraints.

This is an exploratory \emph{single-run} study rather than a multi-seed sweep. All quantitative conclusions below should therefore be interpreted as audited point estimates for one archived run, not as uncertainty-calibrated means.

\section{Results}
\subsection{Main comparison}
\Cref{tab:main-results} reports the full comparison across datasets and steering variants. Three patterns stand out immediately.

First, fixed steering yields substantially larger toxicity reduction than automatic selection. On Civil Comments, fixed global and fixed local reduce toxicity by $0.1923$ and $0.1836$, whereas the corresponding automatic policies achieve $0.0927$ and $0.1093$. On ParaDetox, the gap is even larger: fixed global reaches $0.5139$ and fixed local $0.5850$, while the automatic variants reach $0.2775$ and $0.3332$.

Second, this stronger detox comes with lower semantic preservation. The fixed policies consistently sit lower on the similarity axis than the automatic ones. This is the expected control--preservation trade-off.

Third, locality does not emerge as the decisive improvement. Local steering is sometimes better than global steering, especially on ParaDetox, but the effect is modest relative to the much larger difference between fixed and automatic selection.

\begin{table*}[t]
\centering
\caption{\textbf{Main text detox comparison.} Higher toxicity drop and higher semantic similarity are better. Acceptance rate is the fraction of selected outputs that satisfy the notebook's feasibility constraints.}
\label{tab:main-results}
\small
\begin{tabular}{lllrrrr}
\toprule
Dataset & Direction & Policy & Tox.\ drop & Sem.\ sim. & Accept rate & Mean $\alpha$ \\
\midrule
Civil Comments & Global & Auto & 0.0927 & 0.9563 & 0.0367 & 0.0000 \\
Civil Comments & Global & Fixed & 0.1923 & 0.8912 & 0.0367 & 1.0000 \\
Civil Comments & Local & Auto & 0.1093 & 0.9498 & 0.0333 & 0.0017 \\
Civil Comments & Local & Fixed & 0.1836 & 0.9004 & 0.0333 & 1.0000 \\
ParaDetox & Global & Auto & 0.2775 & 0.8731 & 0.0133 & 0.0000 \\
ParaDetox & Global & Fixed & 0.5139 & 0.7244 & 0.0133 & 1.0000 \\
ParaDetox & Local & Auto & 0.3332 & 0.8512 & 0.0267 & 0.0133 \\
ParaDetox & Local & Fixed & 0.5850 & 0.6960 & 0.0267 & 1.0000 \\
\bottomrule
\end{tabular}
\end{table*}

\begin{figure}[t]
  \centering
  \includegraphics[width=\columnwidth]{figures/tradeoff_toxicity_vs_similarity.png}
  \caption{\textbf{Text detox trade-off: toxicity reduction versus semantic preservation.} Fixed steering moves upward in toxicity reduction but leftward in semantic similarity. Automatic selection is more conservative, especially on Civil Comments.}
  \label{fig:tradeoff}
\end{figure}

\subsection{Automatic step selection collapses toward zero}
The most important diagnostic appears in the selected $\alpha^\star$ values. Global automatic selection chooses $\alpha^\star=0$ for every evaluated example in the archived run. Local automatic selection chooses $\alpha^\star=0$ for $445$ out of $450$ examples and $\alpha^\star=0.5$ only five times. The mean selected step is therefore $0.0000$ for the global policy and $0.0056$ over all local automatic selections, with dataset-specific means of $0.0017$ on Civil Comments and $0.0133$ on ParaDetox.

This behavior is visible in \cref{fig:alpha-hist}. Automatic step selection is not discovering a broad spectrum of safe edit magnitudes. It is overwhelmingly selecting the smallest available step, which in practice means that the system often prefers either the original text or a rewrite generated from an effectively unsteered prompt. That is why the automatic policies score well on similarity but fail to produce strong detox.

\begin{table}[t]
\centering
\caption{\textbf{Automatic-policy collapse diagnostics.} Rates are computed over all selected automatic outputs from both evaluation sets combined.}
\label{tab:collapse}
\small
\begin{tabular}{lrrrr}
\toprule
Mode & $P(\alpha^\star=0)$ & $P(\alpha^\star>0)$ & Rewriter sel. & Original sel. \\
\midrule
Global & 1.0000 & 0.0000 & 0.1800 & 0.8200 \\
Local & 0.9889 & 0.0111 & 0.2156 & 0.7844 \\
\bottomrule
\end{tabular}
\end{table}

\begin{figure}[t]
  \centering
  \includegraphics[width=\columnwidth]{figures/alpha_star_hist.png}
  \caption{\textbf{Automatic $\alpha^\star$ choices.} The automatic policy is overwhelmingly concentrated at zero, especially for the global direction.}
  \label{fig:alpha-hist}
\end{figure}

\subsection{Acceptance rates remain very low}
The feasibility filter is stringent in the current pipeline. Acceptance rates are only $3.67\%$ and $3.33\%$ on Civil Comments, and $1.33\%$ and $2.67\%$ on ParaDetox. Moreover, automatic selection does not improve this quantity over fixed selection in the archived run. This is a useful negative result: the main issue is not just \enquote{pick better alphas}; it is that the overall candidate-generation stage rarely produces outputs that jointly satisfy detox and semantic constraints.

\subsection{Local directions are stable, but not decisive}
The local clustering stage itself appears numerically healthy. \Cref{tab:clusters} summarizes the archived cluster statistics: cluster sizes range from $881$ to $1702$, the local probe AUCs lie between $0.8745$ and $0.9482$, and no cluster falls back to the global direction. So the local method is not failing because of empty or degenerate clusters.

However, this stability does not translate into a decisive end-to-end win. Local fixed steering is slightly worse than global fixed on Civil Comments but better on ParaDetox. Local automatic is slightly better than global automatic on both datasets, but the differences are small compared with the broader problem that automatic selection barely steers at all.

\begin{table}[t]
\centering
\caption{\textbf{Local cluster summary for the archived run.}}
\label{tab:clusters}
\small
\begin{tabular}{lrrrr}
\toprule
$K$ & Min size & Max size & Probe AUC range & Fallback rate \\
\midrule
8 & 881 & 1702 & 0.8745--0.9482 & 0.000 \\
\bottomrule
\end{tabular}
\end{table}

\begin{figure}[t]
  \centering
  \includegraphics[width=\columnwidth]{figures/cluster_sizes.png}
  \caption{\textbf{Cluster sizes for local detox directions.} The local partition is reasonably balanced and does not show pathological tiny clusters.}
  \label{fig:clusters}
\end{figure}

\section{Discussion}
The strongest conclusion of this case study is not that local steering clearly beats global steering. The stronger and more scientifically useful conclusion is that the present automatic policy is too conservative to realize the potential of representation steering. In other words, the main bottleneck lies in the \emph{selection-and-generation loop}, not only in the shape of the direction family.

Why does this happen? The current objective rewards toxicity reduction and semantic preservation simultaneously, but the feasibility thresholds and candidate quality interact in a way that often favors near-identity outputs. If most rewritten candidates are either semantically weak or stylistically unstable, the objective has a strong incentive to stay close to the original sentence. That is exactly what the selected-$\alpha$ histogram shows.

The experiment is still useful because it isolates where the pipeline breaks:
\begin{enumerate}[leftmargin=1.2em,itemsep=0.12em]
\item the representation contains a learnable toxicity signal (probe AUC $0.8931$);
\item fixed steering can exploit that signal to reduce toxicity;
\item but automatic steering collapses because the candidate bank and feasibility rule do not support reliable nontrivial edits.
\end{enumerate}

This interpretation is more informative than a superficial \enquote{the method failed}. It says the method contains a usable control signal, but the current downstream realization of that signal is weak.

\subsection{Threats to validity}
Three caveats matter.

First, this is a single archived run rather than a multi-seed study. The paper therefore supports diagnosis and method refinement more than final quantitative ranking.

Second, the experiment mixes three distinct components: representation steering, a text rewriter, and an external toxicity scorer. A weak result can come from any one of these layers or from their interaction.

Third, the Civil Comments evaluation subset is defined by dataset-side toxicity thresholds, not by the toxicity scorer itself. If the scorer and dataset labels disagree, then the measured \enquote{toxicity drop} can partially reflect a dataset--scorer mismatch rather than pure detox difficulty.

\section{Next steps}
The next iteration should tighten the study scientifically rather than only add more components. The most important next actions are:
\begin{enumerate}[leftmargin=1.2em,itemsep=0.12em]
\item \textbf{Measure nontrivial edit rate explicitly.} Report the fraction of outputs with $\alpha>0$ and the fraction whose selected text differs meaningfully from the original.
\item \textbf{Penalize no-op solutions in the automatic policy.} For example, require a minimal detox gain before allowing $\alpha=0$ to win, or rank feasible nontrivial candidates ahead of identity-like ones.
\item \textbf{Strengthen candidate generation.} The present rewriter frequently produces weak or degraded paraphrases. A stronger rewriting model or better prompts may improve the feasible candidate pool more than a more complicated steering rule.
\item \textbf{Align evaluation subsets with the scorer.} Build a scorer-filtered toxic evaluation split so that the detox metric is evaluated on examples the scorer itself regards as toxic.
\item \textbf{Run multi-seed comparisons.} Once the candidate-generation bottleneck is reduced, rerun the global/local and fixed/auto comparison under multiple seeds.
\end{enumerate}

\section{Conclusion}
This text detox case study shows that representation steering for language is promising but currently unstable. Fixed steering can reduce toxicity, especially on ParaDetox, but pays a semantic cost. Automatic per-sample step selection, which should in principle give a better control--preservation trade-off, instead collapses toward $\alpha^\star\approx 0$ and produces only tiny effective edits. Local directions are numerically well behaved and sometimes slightly better than global ones, but they are not the main story.

The clearest takeaway is therefore methodological: in the current setup, the dominant obstacle is not the absence of locality alone. The dominant obstacle is that the candidate-generation and selection pipeline makes near-identity outputs too attractive. A stronger text steering system will need better candidate generation, better nontrivial-selection logic, and cleaner evaluation alignment.

\section*{Appendix: archived configuration and provenance}
The paper is based on the archived run \texttt{20260405\_205719\_week12\_text\_detox\_steering}. The main configuration choices are:
\begin{itemize}[leftmargin=1.2em,itemsep=0.08em]
\item sentence encoder: \texttt{all-mpnet-base-v2};
\item toxicity scorer: \texttt{unitary/toxic-bert};
\item rewriter: \texttt{google/flan-t5-base};
\item local cluster count: $K=8$;
\item fixed baseline: $\alpha=1.0$;
\item alpha grid: $\{0,0.5,1.0,1.5,2.0\}$;
\item feasibility thresholds: semantic similarity $\ge 0.72$ and toxicity drop $\ge 0.12$;
\item generation: four return sequences and four beams.
\end{itemize}
All scalar values in this report are transcribed from the archived CSV outputs and memo files produced by that run.

\end{document}
